\section{Five truly open questions}

\subsection*{Q1: LET dependence of the direct-damage fraction}
\textbf{Basis.} Paper reports a single 14.2\% mean error against total
SSB/DSB yields across seven LET points but never breaks the total into
direct vs indirect. This is the model's central selling point and the
quantity that discriminates it from empirically tuned predecessors.
\textbf{Next steps.} Request the authors' Geant4-DNA source (data-availability
route) or reimplement the CG geometry from Fig.~S1 in an open TOPAS-nBio
or Geant4-DNA v11 build; score direct vs indirect strand-break
contributions separately at each LET; overlay against TRAX-CHEM and
KURBUC published direct fractions (Nikjoo 2016; Friedland 2017).
Estimated cost: 1--2 GPU-weeks on uicgpu; free per LUCID compute rule.

\subsection*{Q2: Scavenger-concentration modulation}
\textbf{Basis.} Paper validates only against unmodulated dried pBR322 in
air. The claim of correctly separating direct from indirect damage is
untested against the standard scavenger-modulation experiment
(Roots \& Okada 1972; Frankenberg-Schwager 1990; Hirayama 2009). A
genuine direct-only channel must be invariant to scavenger concentration.
\textbf{Next steps.} Add \texttt{G4EmDNAChemistry} scavenger species with
configurable [DMSO] and [glycerol]; re-run the seven LET points at
$\{0, 0.01, 0.1, 1.0\}$~M scavenger; the direct-only lesion count should
be flat within Poisson noise while the total should drop by 2--5$\times$.
If the model's direct channel varies with scavenger concentration, the
direct/indirect partition is not what the paper claims. Free on
Geant4-DNA v11.2 + PolyMorpho chemistry.

\subsection*{Q3: OER prediction across the LET points}
\textbf{Basis.} OER is the strongest observable that discriminates
damage models because it depends on the relative direct vs indirect
contribution \emph{and} on the fixation chemistry. A model that gets
totals right by construction (fit to totals) will not automatically get
OER right unless the direct/indirect split is correctly \emph{derived}.
\textbf{Next steps.} Extend the Geant4-DNA chemistry stage with
anoxic vs aerobic O$_2$ species; compute SSB and DSB yields at 21\%
O$_2$ and $<0.5\%$ O$_2$ for each LET point; compare predicted
OER(LET) to Wenzl \& Wilkens (2011) compilation. This is a
discriminating test the paper never runs.

\subsection*{Q4: Sensitivity to chromatin/plasmid geometry}
\textbf{Basis.} The paper's 90\%-supercoiled + 10\%-relaxed 436-bp
segments in a \SI{3}{\micro\meter} sphere at
$\rho = \SI{2.21e-4}{bp/nm^3}$ is a modelling convenience, not a
physiologically motivated packing. Whether alternative geometries
preserve the 14.2\% agreement is unknown.
\textbf{Next steps.} Re-run the MC with (a) a random-coil pBR322 from a
real solution-state MD trajectory, (b) the DNAFabric random-walk
4361-bp full plasmid without segmentation, (c) a nucleosome-array
analogue as a null; report the 14.2\% metric at each geometry.
If the metric moves $>5\%$, the geometry is doing hidden work.

\subsection*{Q5: Cellular DSB-yield translation}
\textbf{Basis.} Yields are reported per-Gbp on isolated plasmids. The
cellular endpoint (DSB/Gy/cell) is what radiobiologists care about, and
CG-plasmid damage models historically diverge from measured
$\gamma$-H2AX / PFGE yields by 2--4$\times$. The paper never crosses
this bridge.
\textbf{Next steps.} Wrap the CG-plasmid geometry in a nucleus-scale
($5~\mu$m radius) chromatin container at physiological bp density
($\sim\SI{1.5e-2}{bp/nm^3}$, $100\times$ denser than the paper's setup);
rerun at 1~Gy photon exposure; compare predicted DSB/cell against
L\"obrich (2010) $35\pm5$ DSB/Gy/cell benchmark. If the ratio is off by
$>2\times$, the CG model does not translate to the cellular endpoint
and its clinical utility is limited to relative rather than absolute
predictions.
